\section{Logical framework: what a search result means}
\label{sec:framework}

Every condition and every output of the engine carries one of six logical
statuses, and they are never mixed.

\begin{definition}[Statuses of conditions]
A condition used by the search is
\begin{itemize}[leftmargin=*]
\item \st{EXACT} if it is equivalent to the property it tests;
\item \st{NECESSARY} if every valid model satisfies it (so failing it proves
  invalidity, but passing it proves nothing);
\item \st{SUFFICIENT} if satisfying it proves the property (e.g.\ an exact
  rational or real-algebraic slope witness);
\item \st{NUMERIC\_WITNESS} if it is supported only by floating-point evidence;
\item \st{HEURISTIC} if it has no proof in either direction;
\item \st{UNKNOWN} if it could not be decided.
\end{itemize}
Only \st{EXACT} and \st{NECESSARY} conditions may prune. Only \st{EXACT} and
\st{SUFFICIENT} conditions may certify acceptance. \st{HEURISTIC} and
\st{NUMERIC\_WITNESS} conditions may do neither.
\end{definition}

\begin{definition}[Search outcomes]
For a declared domain $\mathcal D$ (geometry, $|\Gam|$, target, Kähler chamber,
and either a box $K$ or a slice $p^\perp$):
\begin{itemize}[leftmargin=*]
\item \st{UNSAT\_BOUNDED}: every state of the finite domain was accounted for and
  none satisfies the declared necessary conditions;
\item \st{UNKNOWN\_BOUNDED}: the domain is exhausted and necessary-condition
  survivors remain, none certified;
\item \st{SAT\_BOUNDED}: a survivor is certified by sufficient conditions
  (never emitted: no model reaches that standard here);
\item \st{UNSAT\_SLICE\_UNBOUNDED}: an entire slice $p^\perp$ has been exhausted
  over \emph{all} integer charges, with no survivor (\cref{sec:slices});
\item \st{UNSAT\_UNBOUNDED}: no model exists for any charges (only the
  divisibility obstruction, \cref{prop:divisibility}, emits this).
\end{itemize}
Stability of a survivor set as $K$ grows is recorded as
\st{EMPIRICAL\_SATURATION} and never converted into an unbounded statement.
\end{definition}

The following observation is elementary but it is what makes the whole design
safe, and it answers precisely when an incomplete test may be used.

\begin{lemma}[Soundness of pruning by necessary conditions]\label{lem:soundness}
Let $\mathcal D$ be finite and let $R_1,\dots,R_n$ be conditions each of which is
necessary for the declared target under stated assumptions. If the search
enumerates $\mathcal D$ and discards a state only when some $R_i$ fails---on the
state or on a prefix from which the failure provably propagates to every
completion---then every valid model in $\mathcal D$ survives. In particular
\st{UNSAT\_BOUNDED} is correct. Replacing any $R_i$ by a weaker necessary
condition can only enlarge the survivor set.
\end{lemma}

\begin{proof}
A valid model satisfies every $R_i$, so it is never discarded at the leaf. A
prefix is discarded only if every completion fails some $R_i$, which a valid
model does not. Weakening a necessary condition removes rejections and adds
none.
\end{proof}

The last sentence is the answer to a question that matters in practice: an
\emph{incomplete} test (one that misses some obstructions) is harmless for
\st{UNSAT\_BOUNDED}; only an \emph{unsound} test (one that rejects a state that
could be valid) is dangerous. The finite coefficient box of the slope test
(\cref{sec:slope}) is incomplete but sound, and may therefore contribute to
\st{UNSAT\_BOUNDED}; its completion is nevertheless worth having, because an
incomplete test leaves false survivors, and \cref{sec:validation} shows that it
does.
